\documentclass[10pt,a4paper]{amsart}
\usepackage[margin=19mm]{geometry}
\usepackage{amsmath,amssymb,mathtools}
\usepackage[scr=rsfs,cal=euler]{mathalfa}
\usepackage{xfrac}
\usepackage[unicode,hidelinks,bookmarks=false]{hyperref}
\newcommand{\ie}{i.e.\ }
\newcommand{\cf}{cf.\ }
\newcommand{\resp}{respectively\ }
\newcommand{\Subsection}{\subsection}
\providecommand{\nobreakdash}{\nobreak\hbox{-}}
\setlength{\emergencystretch}{2em}
\multlinegap0pt
\newcommand{\skpt}{\par\smallskip}
\providecommand{\eg}{e.g.\ }
\usepackage{enumerate}
\usepackage{pgfplots}
\usepackage{mathrsfs}
\usepackage{xcolor}
\usepackage{tikz-cd}
\usepackage{scalerel}
\usepackage{circuitikz}

%%%%%%%%%%%%%%%%%%%%%%%%%%%%%%%%%%%%%%%%%
\definecolor{lblue}{rgb}{0.0, 0.55, 1.0}
\definecolor{dred}{rgb}{1.0, 0.0, 0.3}
\definecolor{byzant}{rgb}{0.74, 0.2, 0.64}
\definecolor{pomme}{rgb}{0.16, 0.67, 0.53}
\definecolor{aqua}{rgb}{0.0, 1.0, 1.0}
\definecolor{azure}{rgb}{0.0, 0.5, 1.0}
\definecolor{carnelian}{rgb}{0.7, 0.11, 0.11}
\definecolor{ceruleanblue}{rgb}{0.16, 0.32, 0.75}
\definecolor{pgreen}{rgb}{0.0, 0.65, 0.58}
\definecolor{pblue}{rgb}{0.11, 0.22, 0.73}
\definecolor{pear}{rgb}{0.82, 0.89, 0.19}
\definecolor{pyellow}{rgb}{1.0, 0.97, 0.12}
\definecolor{lgreen}{rgb}{0.75, 1.0, 0.0}
\definecolor{brown}{rgb}{0.6, 0.4, 0.08}
\definecolor{ggray}{rgb}{0.27, 0.35, 0.27}

\definecolor{lblue}{rgb}{0.0, 0.55, 1.0}
\definecolor{dred}{rgb}{1.0, 0.0, 0.3}
\definecolor{byzant}{rgb}{0.74, 0.2, 0.64}
\definecolor{pomme}{rgb}{0.16, 0.67, 0.53}
\definecolor{aqua}{rgb}{0.0, 1.0, 1.0}
\definecolor{azure}{rgb}{0.0, 0.5, 1.0}
\definecolor{carnelian}{rgb}{0.7, 0.11, 0.11}
\definecolor{ceruleanblue}{rgb}{0.16, 0.32, 0.75}
\definecolor{pgreen}{rgb}{0.0, 0.65, 0.58}
\definecolor{pblue}{rgb}{0.11, 0.22, 0.73}
\definecolor{pear}{rgb}{0.82, 0.89, 0.19}
\definecolor{pyellow}{rgb}{1.0, 0.97, 0.12}
\definecolor{lgreen}{rgb}{0.75, 1.0, 0.0}
\definecolor{brown}{rgb}{0.6, 0.4, 0.08}
\definecolor{ggray}{rgb}{0.27, 0.35, 0.27}
\definecolor{darkgreen}{HTML}{00AA00}

%%%%%%%%%%%%%%%%%%%%%%%%%%%%%%%%%%%%%%%%

\DeclareRobustCommand{\SkipTocEntry}[5]{}


\ctikzset{bipoles/resistor/height=0.15}
\ctikzset{bipoles/resistor/width=0.4}
\usetikzlibrary {positioning}
\usetikzlibrary{patterns}
\pgfplotsset{compat=1.10}
\usepgfplotslibrary{fillbetween}

\usetikzlibrary{patterns}
\usetikzlibrary{calc,decorations.markings}
\usetikzlibrary{shapes,snakes}


\tikzstyle{Cwhite}=[scale =.8,circle, fill = white, minimum size=3mm]
\tikzstyle{Cgray}=[scale =.4,circle, fill = gray, minimum size=3mm]
\tikzstyle{Cblack2}=[scale =.4,circle, fill = black, minimum size=5mm]
\tikzstyle{Cblack}=[scale =.7,circle, fill = black, minimum size=3mm]
\tikzstyle{C0}=[scale =.9,circle, fill = black!0, inner sep = 0pt, minimum size=3mm]
\tikzstyle{C1}=[scale =.7,circle, fill = black!0, inner sep = 0pt, minimum size=3mm]
\tikzstyle{Cred}=[scale =.4,circle, fill = red, minimum size=3mm]


\newcommand{\symbuptm}{\tikz[scale=.37, baseline=-0.7]{\draw(.1,-.05)to[out=50,in=-130](.9,.05) (.5,0)--++(0,.5);}}
%\newcommand{\symbuptm}{\tikz[scale=.28, baseline=-0.7]{\draw(.1,-.05)to[out=50,in=-130](.9,.05) (.5,0)--++(0,.5);}}
%\newcommand{\almorth}{\tikz[scale=.28, baseline=-0.7]{\symbuptm \draw (-.44,.22) circle (14pt);}}
\newcommand{\almorth}{\tikz[scale=.37, baseline=-0.7]{\symbuptm \draw (-.44,.22) circle (14pt);}}
\newcommand{\aplus}{\hspace{-.2cm}\almorth\,}
\newcommand{\aperp}{\,\symbuptm\,}

%%%%%%%%%%%%%%%%%%%%%%%%%%%%%%%%%%%%%%%%%%%%%%%%%%%%%%%

%%%%%%%%%%%%%%%%%%%%%%%%%%%%%%%%%%%%%%%%%%%%%%%%%%%%%%%%%%\langleetc modifiés
\newcommand{\innone}[2]{{\langle#2\rangle_{#1}}}

\newcommand{\binnone}[2]{{\left\langle#2\right\rangle_{#1}}}

\newcommand{\llangle}{\mathopen{\langle\mkern -3.1mu\lvert}}
\newcommand{\rrangle}{\mathclose{\rvert\mkern -3.1mu\rangle}}

\newcommand{\bllangle}{\mathopen{\bigl\langle\mkern -3.1mu\bigl\lvert}}
\newcommand{\brrangle}{\mathclose{\bigr\rvert\mkern -3.1mu\bigr\rangle}}

\newcommand{\highinn}[2]{{\llangle#2\rrangle_{#1}}}
\newcommand{\bhighinn}[2]{{\bllangle#2\brrangle_{#1}}}
\newcommand{\highinnoneorth}[2]{{\llangle#2\rrangle^{\perp}_{#1}}}
\newcommand{\innoneorth}[2]{{\langle#2\rangle^{\perp}_{#1}}}



\newcommand{\symbupl}{\tikz[scale=.1, baseline=-1.7]{\draw(.98,-.75)to[out=90,in=90](.98,.86) ;}}
\newcommand{\langleup}{\langle \tikz[scale=.28, baseline=-0.7]{\symbupl}}
\newcommand{\symbupr}{\tikz[scale=.1, baseline=-1.7]{\draw(.91,-.77)to[out=90,in=90](.91,.86) ;}}
\newcommand{\rangleup}{\rangle\tikz[scale=.28, baseline=-0.7]{\symbupr}}

\newcommand{\highinnprime}[2]{{\llangle #2\rrangle'_{#1}}}
%%%%%%%%%%%%%%%%%%%%%%%%%%%%%%%%%%%%%%%%%%%%%%%MODIF DE \bar pour overbar
\makeatletter
\newsavebox\myboxA
\newsavebox\myboxB
\newlength\mylenA

\newcommand*\overbar[2][0.75]{
\sbox{\myboxA}{$\m@th#2$}
\setbox\myboxB\null
\ht\myboxB=\ht\myboxA
\dp\myboxB=\dp\myboxA
\wd\myboxB=#1\wd\myboxA
\sbox\myboxB{$\m@th\overline{\copy\myboxB}$}
\setlength\mylenA{\the\wd\myboxA}
\addtolength\mylenA{-\the\wd\myboxB}
\ifdim\wd\myboxB<\wd\myboxA
\rlap{\hskip 1\mylenA\usebox\myboxB}{\usebox\myboxA}
\else
\hskip -0.5\mylenA\rlap{\usebox\myboxA}{\hskip 0.5\mylenA\usebox\myboxB}
\fi}
\makeatother
\newcommand{\grind}[1]{#1}
\newcommand{\comp}[1]{\overbar[.5]{#1}}
%%%%%
\newcommand{\Hyb}{{\mathit{hyb}}}
\newcommand{\hyb}{^{\Hyb}}

\newcommand{\mg}{\mathscr M}
\newcommand{\mgbar}{\comp{\mathscr M}}
\newcommand{\mgg}[1]{\mg_{#1}^{ }}
\newcommand{\mgbarg}[1]{\mgbar_{#1}^{ }}
\newcommand{\mghyb}[1]{\mg_{#1}^{\hyb}}

%%%%%%%%%%%%%%%%%%%%%%%%%%%%%%%%%%%%%%%%%%%%%%%%%%%%%%%
\newcommand{\norm}[1]{\|#1\|}
\newcommand{\normsq}[1]{\|#1\|^{2}}
\newcommand{\normbsq}[1]{\left\|#1\right\|^{2}}
\newcommand{\refnorm}[1]{\lvert#1\rvert}
\newcommand{\refnormsq}[1]{\lvert#1\rvert^{2}}
\newcommand{\refnormbsq}[1]{\left\lvert#1\right\rvert^{2}}
\newcommand{\normorth}[1]{\|#1\|^{\perp}}

\newcommand{\rest}[1]{\vert_{#1}}

\newcommand{\abs}[1]{\lvert #1\rvert}
\newcommand{\abss}[1]{\left\lvert #1\right\rvert}



\newcommand{\rquot}[2]{#1/\,#2}
\newcommand{\st}{\bigm|}

\newcommand{\orient}{\mathcal{O}}
%%%%%%%%%%%%%%%%%%%%%%%%%%%%%%%%%%%%%%%%%%%%
\DeclareMathOperator{\rk}{rank}
\DeclareMathOperator{\Hom}{Hom}
\DeclareMathOperator{\id}{Id}
\DeclareMathOperator{\weak}{w}
\DeclareMathOperator{\projeuc}{p}
\DeclareMathOperator{\Jac}{Jac}
\DeclareMathOperator{\vol}{vol}
\DeclareMathOperator{\Vor}{Vor}

\DeclareMathOperator{\spann}{span}

\renewcommand{\Re}{\operatorname{Re}}
\DeclareMathOperator{\intt}{int}
\DeclareMathOperator{\dist}{d}
\newcommand{\graphgenus}{{h}}
\newcommand{\filter}{\mathrm{F}}
\newcommand{\Q}{\mathbb{Q}}
\newcommand{\R}{\mathbb{R}}
\newcommand{\Z}{\mathbb{Z}}
\newcommand{\N}{\mathbb{N}}
\newcommand{\C}{\mathbb{C}}
\newcommand{\E}{\mathbb{E}}
\newcommand{\T}{\mathbb{T}}
\renewcommand{\L}{\mathbb{L}}
\renewcommand{\S}{\mathbb{S}}

\newcommand{\smallcc}{\mathbb{C}}
\newcommand{\mgr}{\mathcal{G}}
\newcommand{\curve}{\mathcal{C}}
\newcommand{\rsf}{\mathcal{S}}
\newcommand{\Ical}{\mathcal{I}}

\newcommand{\grm}[2]{\mathrm{gr}_{#1}^{#2}}


\newcommand{\proj}{\mathfrak{p}}
\newcommand{\qf}{\mathfrak{q}}
\newcommand{\fin}{\mathfrak f}
\newcommand{\scS}{\mathscr{S}}

\newcommand{\transpose}{\mathrm{T}}

\renewcommand{\setminus}{\smallsetminus}
\renewcommand{\emptyset}{\varnothing}
\newcommand{\dtor}{\zeta}

\newcommand{\interior}[1]{\intt(#1)}






\newcommand{\hdist}{d_{\mathrm{H}}}
\newcommand{\GH}{\mathrm{GH}}
\newcommand{\HA}{\mathrm{H}}



\newcommand{\suppp}[1]{{\mathrm{supp}(#1)^+}}
\newcommand{\suppm}[1]{{\mathrm{supp}(#1)^-}}
\newcommand{\supppm}[1]{{\mathrm{supp}(#1)^\pm}}



\newcommand{\lexeq}{\preceq_{\mathrm{lex}}}
\newcommand{\lexst}{\prec_{\mathrm{lex}}}
\newcommand{\gexeq}{\succeq_{\mathrm{lex}}}
\newcommand{\gex}{\succ_{\mathrm{lex}}}



%%%%%%%%%%%%%%%%%%%%%%%%%%%%%%%%%%%%%%%%%%%%%%%%%
\renewcommand*{\to}{\mathchoice{\longrightarrow}{\rightarrow}{\rightarrow}{\rightarrow}}
%%%%%%%%%%
\renewcommand*{\implies}{\mathchoice{\Longrightarrow}{\Rightarrow}{\Rightarrow}{\Rightarrow}}
%%%%%%%%%%
\let\oldmapsto\mapsto
\renewcommand*{\mapsto}{\mathchoice{\longmapsto}{\oldmapsto}{\oldmapsto}{\oldmapsto}}
%%%%%%%%%%%
\newcommand{\longhookrightarrow}{\lhook\joinrel\longrightarrow}
\newcommand*{\hookto}{\mathchoice{\longhookrightarrow}{\hookrightarrow}{\hookrightarrow}{\hookrightarrow}}
%%%%%
%Better widebar

\makeatletter
\let\save@mathaccent\mathaccent
\newcommand*\if@single[3]{%
  \setbox0\hbox{${\mathaccent"0362{#1}}^H$}%
  \setbox2\hbox{${\mathaccent"0362{\kern0pt#1}}^H$}%
  \ifdim\ht0=\ht2 #3\else #2\fi
  }
%The bar will be moved to the right by a half of \macc@kerna, which is computed by amsmath:
\newcommand*\rel@kern[1]{\kern#1\dimexpr\macc@kerna}
%If there's a superscript following the bar, then no negative kern may follow the bar;
%an additional {} makes sure that the superscript is high enough in this case:
\newcommand*\widebar[1]{\@ifnextchar^{{\wide@bar{#1}{0}}}{\wide@bar{#1}{1}}}
%Use a separate algorithm for single symbols:
\newcommand*\wide@bar[2]{\if@single{#1}{\wide@bar@{#1}{#2}{1}}{\wide@bar@{#1}{#2}{2}}}
\newcommand*\wide@bar@[3]{%
  \begingroup
  \def\mathaccent##1##2{%
%Enable nesting of accents:
    \let\mathaccent\save@mathaccent
%If there's more than a single symbol, use the first character instead (see below):
    \if#32 \let\macc@nucleus\first@char \fi
%Determine the italic correction:
    \setbox\z@\hbox{$\macc@style{\macc@nucleus}_{}$}%
    \setbox\tw@\hbox{$\macc@style{\macc@nucleus}{}_{}$}%
    \dimen@\wd\tw@
    \advance\dimen@-\wd\z@
%Now \dimen@ is the italic correction of the symbol.
    \divide\dimen@ 3
    \@tempdima\wd\tw@
    \advance\@tempdima-\scriptspace
%Now \@tempdima is the width of the symbol.
    \divide\@tempdima 10
    \advance\dimen@-\@tempdima
%Now \dimen@ = (italic correction / 3) - (Breite / 10)
    \ifdim\dimen@>\z@ \dimen@0pt\fi
%The bar will be shortened in the case \dimen@<0 !
    \rel@kern{0.6}\kern-\dimen@
    \if#31
      \overline{\rel@kern{-0.6}\kern\dimen@\macc@nucleus\rel@kern{0.4}\kern\dimen@}%
      \advance\dimen@0.4\dimexpr\macc@kerna
%Place the combined final kern (-\dimen@) if it is >0 or if a superscript follows:
      \let\final@kern#2%
      \ifdim\dimen@<\z@ \let\final@kern1\fi
      \if\final@kern1 \kern-\dimen@\fi
    \else
      \overline{\rel@kern{-0.6}\kern\dimen@#1}%
    \fi
  }%
  \macc@depth\@ne
  \let\math@bgroup\@empty \let\math@egroup\macc@set@skewchar
  \mathsurround\z@ \frozen@everymath{\mathgroup\macc@group\relax}%
  \macc@set@skewchar\relax
  \let\mathaccentV\macc@nested@a
%The following initialises \macc@kerna and calls \mathaccent:
  \if#31
    \macc@nested@a\relax111{#1}%
  \else
%If the argument consists of more than one symbol, and if the first token is
%a letter, use that letter for the computations:
    \def\gobble@till@marker##1\endmarker{}%
    \futurelet\first@char\gobble@till@marker#1\endmarker
    \ifcat\noexpand\first@char A\else
      \def\first@char{}%
    \fi
    \macc@nested@a\relax111{\first@char}%
  \fi
  \endgroup
}
\makeatother

\let\oldbar\bar
\renewcommand*{\bar}[1]{{\mathchoice{\widebar{#1}}{\widebar{#1}}{\widebar{#1}}{\oldbar{#1}}}}

\newcommand*{\mk}{\mkern -1mu}
\newcommand*{\Mk}{\mkern -2mu}
\newcommand*{\mK}{\mkern 1mu}
\newcommand*{\MK}{\mkern 2mu}

\newcommand*{\relabel}{\renewcommand{\labelenumi}{(\theenumi)}}
\newcommand*{\romanenumi}{\renewcommand*{\theenumi}{\roman{enumi}}\relabel}
\newcommand*{\Romanenumi}{\renewcommand*{\theenumi}{\Roman{enumi}}\relabel}
\newcommand*{\alphenumi}{\renewcommand*{\theenumi}{\alph{enumi}}\relabel}
\newcommand*{\alphebisnumi}{\renewcommand*{\theenumi}{\alph{enumi}$'$}\relabel}
\newcommand*{\Alphenumi}{\renewcommand*{\theenumi}{\Alph{enumi}}\relabel}
\let\oldtilde\tilde
\renewcommand*{\tilde}[1]{\mathchoice{\widetilde{#1}}{\widetilde{#1}}{\oldtilde{#1}}{\oldtilde{#1}}}
\let\oldhat\hat
\renewcommand*{\hat}[1]{\mathchoice{\widehat{#1}}{\widehat{#1}}{\oldhat{#1}}{\oldhat{#1}}}
\let\oldexists\exists
\renewcommand*{\exists}{\mathrel{\oldexists}}
\let\oldforall\forall
\renewcommand*{\forall}{\mathrel{\oldforall}}

\newtheorem{cdrthm}{Theorem}[section]
\newtheorem{theo}[cdrthm]{Theorem}
\newtheorem{lemm}[cdrthm]{Lemma}
\newtheorem{prop}[cdrthm]{Proposition}
\theoremstyle{definition}\newtheorem{defi}[cdrthm]{Definition}
\newtheorem{rema}[cdrthm]{Remark}
\newtheorem{conv}[cdrthm]{Convention}
\title[Finite nearest sites in degenerations]{Higher rank inner products, Voronoi tilings and metric degenerations of tori: original Lemma6.1}
\author{Omid Amini and Noema Nicolussi}
\date{Source: Annales Henri Lebesgue8(2025),1109--1188; DOI10.5802/ahl.257}
\begin{document}\maketitle
\noindent\textit{Editorial scope.} The counted unit is exact original Lemma6.1 for weakly tame families with all original scale, graded discreteness and compact-set quantifiers. The complete canonical-lift and coercivity core is uncounted support. Finite illustrated examples and later Voronoi/Hausdorff/tori applications are omitted. Every general constraint used here is textual; native mathematical symbols and original index/constant notation remain unchanged.
\section{Original basic notation}
\subsection*{Basic notations} \label{ss:BasicNotations}
For a point $a$ in the vector space $\R^r$, we denote by $a_1, \dots, a_r$ the coordinates of~$a$. We set $\R_+ \coloneqq (0, +\infty)$ and denote by $\R_+^r =(0,+\infty)^r$ the set of $r$-vectors with positive coordinates.


We use the notation $\highinn{}{\cdot\,,\cdot}$ for inner products with values in $\R^r$, $r\in \N$. For scalar products, we usually use $\innone{}{\cdot\,,\cdot}$.

\section{Higher rank inner products} \label{sec:HigherRankInnerProducts}

\subsection{Totally ordered vector spaces}



Let $(\Lambda, \preceq)$ be a totally ordered real vector space of finite dimension. It is known that $(\Lambda, \preceq)$ is isomorphic to $\R^r$ with its lexicographic order $\lexeq$ where $r$ is the dimension of $\Lambda$~\cite{Bir-Lattice, HW52}. This means there is no loss of generality in assuming that $\Lambda = \R^r$ and $\preceq =\lexeq$. For $a =(a_1, \dots, a_r)$ and $b=(b_1, \dots, b_r)$ in $\R^r$, we have $a \preceq b$ if either $a=b$, or there exists $j\in [r]$ such that $a_1=b_1, \dots, a_{j-1}=b_{j-1}$ and $a_j<b_j$.



The vector space $\Lambda = \R^r$ with its lexicographic order $\preceq=\lexeq$ has a natural decreasing filtration
\begin{equation}\label{eq:FiltrationRr}
\Lambda^\bullet \colon \, \Lambda^1 \coloneq \Lambda\supset \Lambda^2 \supset \dots \supset \Lambda^{r} \supset \Lambda^{r+1} \coloneq (0)
\end{equation}
defined by
\[
\Lambda^j \coloneqq \bigl\{a = (a_1, \dots, a_r) \in \R^r \bigm| a_1 = \dots =a_{j-1}=0\bigr\}.
\]
Note that $\Lambda^j$ is naturally isomorphic to $\R^{r+1-j}$ endowed with its lexicographic order.



\subsection{Inner product spaces of higher rank}

Let $H$ be a real vector space of finite dimension and $\Lambda =\R^r$ with its lexicographical order $\preceq=\lexeq$.

Consider a symmetric bilinear form $\highinn{}{\cdot \,, \cdot} \colon H \times H \to \R^r$. We denote by $\highinn{j}{\cdot\,, \cdot}$ the $j^{\rm th}$ coordinate of $\highinn{}{\cdot\,,\cdot}$, so that $\highinn{}{\cdot\,,\cdot}$ has the form
\[
\highinn{}{\cdot\,,\cdot} = \bigl(\highinn{1}{\cdot\,, \cdot}, \dots, \highinn{r}{\cdot\,, \cdot}\bigr).
\]
In case that $r=1$, we call $\highinn{}{\cdot \,, \cdot}$ a \emph{scalar form}.

We say that $\highinn{}{\cdot \,, \cdot}$ is \emph{non-negative} if $\highinn{}{x\,, x} \succeq 0$ for all $x \in H$, and \emph{positive} if $\highinn{}{x\,, x} \succ 0$ for all $x \in H \setminus \{0\}$. In the case $r=1$, a positive scalar form is called a \emph{scalar product}.

A symmetric bilinear form $\highinn{}{\cdot \,, \cdot} \colon H \times H \to \R^r$ defines a filtration on $H$ as follows. Consider the decreasing filtration $\Lambda^\bullet$ on $\Lambda = \R^r$ given in~\eqref{eq:FiltrationRr}. Setting
\begin{align}
\filter^j \,\MK &\!\Mk:= \bigl\{ x \in H \bigm| \,\highinn{}{x \,, y} \in \Lambda^j \text{ for all }y \in H\bigr\}\\
&= \bigl\{ x \in H \bigm| \,\highinn{}{x \,, y}_i = 0 \quad\text{for all }i <j\text{  and all }y \in H\bigr\}
\end{align}
we obtain a non-increasing filtration
\[
\filter^\bullet\colon \filter^1 = H \supseteq \filter^2 \supseteq \filter^3 \supseteq \dots \supseteq \filter^r \supset \filter^{r+1} = (0)
\]
of $H$. We call $\filter^\bullet$ the \emph{filtration induced by} $\highinn{}{\cdot \,, \cdot}$.

\setcounter{cdrthm}{3}
The following properties are immediate from the definition of the filtration $\filter^\bullet$.

\begin{prop}\label{prop:GradedForms}
Let $\highinn{}{\cdot \,, \cdot} \colon H \times H \to \R^r$ be a symmetric bilinear form. For $j \in [r]$, the $j^{\rm th}$ component $\highinn{j}{\cdot \,, \cdot}$ of $\highinn{}{\cdot \,, \cdot}$ induces a symmetric bilinear form
\[
\highinn{j}{\cdot\,, \cdot} \colon \rquot{\filter^j}{\filter^{j+1}} \times \rquot{\filter^j}{\filter^{j+1}} \to \R
\]
on the quotient $\rquot{\filter^j}{\filter^{j+1}}$.

If $\highinn{}{\cdot \,, \cdot}$ is non-negative on $H$, then $\highinn{j}{\cdot \,, \cdot}$ is non-negative on $\rquot{\filter^j}{\filter^{j+1}}$ for all $j$. If $\highinn{j}{\cdot \,, \cdot}$ is positive on $\rquot{\filter^j}{\filter^{j+1}}$ for all $j$, then $\highinn{}{\cdot \,, \cdot}$ is positive on $H$.
\end{prop}



\begin{defi}[Inner product of higher rank] \label{defi:InnerProduct}\leavevmode
\begin{itemize}
\item An \emph{inner product on $H$ with values in $\Lambda = \R^r$} is a symmetric bilinear form
\[
\highinn{}{\cdot\,, \cdot} \colon H \times H \to \R^r
\]
such that for all $j \in [r]$, the induced form $\highinn{j}{\cdot\,, \cdot}$ on $\filter^j / \filter^{j+1}$ is positive, equivalently, if $\highinn{j}{\gamma, \gamma} >0$ for all $\gamma \in \filter^j \setminus \filter^{j+1}$. The pair $(H, \highinn{}{\cdot\,, \cdot})$ is called \emph{an inner product space}.

\item For $j \in [r]$, the $j^{\rm th}$ \emph{graded piece} $\grm{}{j}H$ of $(H, \highinn{}{\cdot\,, \cdot})$ is the quotient
\[
\grm{}{j}H = \rquot{\filter^j}{\filter^{j+1}}
\]
endowed with the scalar product
\[
\highinn{j}{\cdot\,, \cdot} \colon \grm{}{j}H\times \grm{}{j}H \to \R.
\]
We denote by $\proj_j \colon \filter^j \to \grm{}{j}H$ the projection map.

\item By a slight abuse of notation, we refer to inner products with values in $\R^r$, for arbitrary positive integer $r$, as \emph{inner products of higher rank} or simply as \emph{inner products} (although, strictly speaking, $r=1$ is allowed and, for larger~$r$, their image is allowed to be a one-dimensional subspace of $\R^r$).%\qedd
\end{itemize}
\end{defi}

\setcounter{cdrthm}{8}
\begin{defi}[Quadratic form associated to an inner product] For an inner product $\highinn{}{\cdot\,, \cdot} \colon H \times H \to \R^r$, we denote by $\qf$ its \emph{quadratic form}
\begin{align*}
\qf = (\qf_1, \dots, \qf_r) \colon H &\to  \R^r \\
\gamma &\mapsto \highinn{}{\gamma,\gamma}.%\qedd
\end{align*}
\end{defi}

\setcounter{cdrthm}{16}
\subsection{The almost orthogonal decomposition} \label{ss:AlmostOrthogonalDecomposition}
Let $\highinn{}{\cdot\,,\cdot} \colon H \times H \to \R^r$ be an inner product and $\filter^\bullet$ the induced filtration on $H$.



Consider two elements $\gamma \in \filter^{j} \setminus \filter^{j+1}$ and $\gamma' \in \filter^{j'} \setminus \filter^{j'+1}$. Then $\highinn{}{\gamma\,,\gamma'}$ lies in ${\Lambda}^{\max\{j, j'\}}$, that is, it has the form
\[
\highinn{}{\gamma\,,\gamma'} = \bigl(0, \dots, 0, \highinn{\max\{j, j'\}}{\gamma\,,\gamma'}, \dots, \highinn{r}{\gamma\,,\gamma'} \bigr).
\]
This motivates the following definition.

\begin{defi}[Almost orthogonality]
Two elements $\gamma, \gamma' \in H$ are called \emph{almost orthogonal} and we write $\gamma \aperp\gamma'$ if $\highinn{\max\{j,j'\}}{\gamma, \gamma'} = 0$ where $j, j' \in [r]$ are the indices with $\gamma \in \filter^j \setminus \filter^{j+1}$ and $\gamma' \in \filter^{j'} \setminus \filter^{j'+1}$.

Two subsets $A, B \subset E$ are called \emph{almost orthogonal} and we write $A \aperp B$, if any pair of elements $\gamma \in A$ and $\eta \in B$ are almost orthogonal. We use the notation $A \aplus B$ for $A+B$, indicating that $A$ and $B$ are almost orthogonal.%\qedd
\end{defi}

The terminology is justified as follows.

\begin{lemm}\label{lem:orthogonality_vanishing}
If $\gamma \in H$ is almost orthogonal to itself, then $\gamma = 0$. In particular, if $U, V \subset H$ are almost orthogonal subspaces, then $U \aplus V$ is a direct sum. That is, if $u+v = 0$ for $u\in U$ and $v \in V$, then $u=v = 0$.
\end{lemm}

\begin{proof}
If $\gamma \neq 0$, then there exists $k \in [r]$ with $\gamma \in \filter^k \smallsetminus \filter^{k+1}$. Moreover, $\highinn{k}{\gamma, \gamma} >0$ since $\highinn{}{\cdot, \cdot}$ is an inner product. On the other hand, if $\gamma$ is orthogonal to itself, then $\highinn{k}{\gamma, \gamma} =0$. We have arrived at a contradiction.
\end{proof}


As we discuss next, the notion of almost orthogonality allows to canonically decompose $H$ in terms of its graded pieces $\grm{}{j}H$, $j \in [r]$.

\begin{lemm}[Lifting lemma] \label{lem:LiftingLemma}
Let $\highinn{}{\cdot\,,\cdot} \colon H \times H \to \R^r$ be an inner product. Then for $j\in [r]$, there exists a unique linear map $\proj_j^* \colon \grm{}{j}H \to \filter^j$ with the following properties:
\begin{enumerate}
\item $\proj_j \circ \proj_j^* = \id$ and, in particular, $\proj_j^* \colon \grm{}{j}H \hookto \filter^j$ is an embedding.
\item The image of $\proj_j^*$ is almost orthogonal to $\filter^{j+1}$.
\end{enumerate}
The inner product space $(H, \highinn{}{\cdot, \cdot})$ has the almost orthogonal decomposition
\[
H = \proj_1^*\left(\grm{}{1}H\right) \aplus\,\proj_2^*\left(\grm{}{2}H\right) \aplus \cdots\!\aplus\, \proj_r^*\left(\grm{}{r}H\right).
\]
That is, $H$ is the direct sum of the subspaces $\proj_j^*(\grm{}{j}(H))$, $j \in [r]$, and they are pairwise almost orthogonal.
\end{lemm}

\begin{defi}
We refer to the embedding $\proj_j^*\colon \grm{}{j} H \hookto H$ as the \emph{canonical lifting} of $\grm{}{j}H$ to $H$, and denote its image by $H_j \coloneqq \proj_j^*(\grm{}{j} H)$. The resulting almost orthogonal decomposition is denoted by
\[
H = H_1\aplus H_2 \aplus \cdots \aplus H_r. %\qedd
\]
\end{defi}

Note that the $j^{\rm th}$ component $\highinn{j}{\cdot \,, \cdot}$ of the inner product $\highinn{}{\cdot \,, \cdot}$ is a scalar product on $H_j$ for all $j\in [r]$. Namely, by the properties in Lemma~\ref{lem:LiftingLemma}, it is clear that $H_j \subset F^j$ and $H_j \aperp F^{j+1}$. Applying Lemma~\ref{lem:orthogonality_vanishing}, it follows that every $x \in H \setminus \{0\}$ belongs to $F^j \setminus F^{j+1}$ and consequently satisfies $\highinn{j}{x \,, x}>0$.


\begin{proof}[Proof of Lemma~\ref{lem:LiftingLemma}]
Fix $j \in [r]$ and let $\gamma\in \grm{}{j}H$. We show that there exists a unique element $\theta \in \filter^j$ such that $\proj_j(\theta) =\gamma$, and $\theta$ is almost orthogonal to $\filter^{j+1}$. This defines a unique linear map $\proj_j^*$ with the above properties.

To prove the uniqueness of $\theta$, let $\theta$ and $\theta'$ be two elements of $\filter^j$ such that $\proj_j(\theta) =\proj_j(\theta') = \gamma$, and both are almost orthogonal to $\filter^{j+1}$. Then the difference $\theta-\theta'$ belongs to $\filter^{j+1}$ and is moreover almost orthogonal to $\filter^{j+1}$. Lemma~\ref{lem:orthogonality_vanishing} then ensures that $\theta =\theta'$.


In order to prove the existence of $\theta$, we proceed inductively, and show that for each $k=0, \dots, r-j$, we can find an element $\theta_{j+k}\in \filter^j$ such that
\begin{enumerate}\romanenumi
\item\label{prflemm2.19.1} $\proj_j(\theta_{j+k}) = \gamma$, and

\item\label{prflemm2.19.2} for each $h=1, \dots, k$, we have $\highinn{j+h}{\theta_{j+k}, \eta} = 0 \text{ for all }\eta\in \filter^{j+h}.$
\end{enumerate}
Setting $\theta = \theta_{r}$, we then obtain the desired element.

For $k=0$, choose any $\theta_j \in \filter^j$ with $\proj_j(\theta_j) =\gamma$. Then $\theta_j$ verifies~\eqref{prflemm2.19.1} and~\eqref{prflemm2.19.2} is empty.

Assume we have found $\theta_{j+k}$ for some $k< r-j$. Define the linear form
\begin{align*}
f_{j+k}\colon \filter^{j+k+1} &\to \R\\
\eta &\longmapsto \highinn{j+k+1}{\theta_{j+k}, \eta}.
\end{align*}
Since $\highinn{j+k+1}{\cdot \,, \cdot}$ is an inner product on $\grm{}{j+k+1}H$, there exists an element $\alpha \in \filter^{j+k+1}$ with
\[
f_{j+k}(\eta) = \highinn{j+k+1}{\alpha, \eta}, \quad \forall \eta\in \filter^{j+k+1}.
\]
Let $\theta_{j+k+1} \coloneqq \theta_{j+k} - \alpha$ and note that $\proj_j(\theta_{j+k+1}) = \proj_j(\theta_{j+k}) =\gamma$. Since $\alpha \in \filter^{j+k+1}$, we also have $\highinn{j+h}{\alpha, \eta}=0$ for any $\eta \in \filter^{j+h}$ for all $h=1, \dots, k$. Finally,
\[
\highinn{j+k+1}{\theta_{j+k+1}, \eta} = \highinn{j+k+1}{\theta_{j+k}, \eta} - \highinn{j+k+1}{\alpha, \eta}= 0, \quad \forall \eta\in \filter^{j+k+1},
\]
by construction. Thus, $\theta_{j+k+1}$ satisfies the properties~\eqref{prflemm2.19.1} and~\eqref{prflemm2.19.2}. This finishes the proof of existence and uniqueness of the embedding $\proj_j^* \colon \grm{}{j}H \hookto H$.

The properties of $\proj_j^*$ clearly imply that the subspaces $\proj_i^*(\grm{}{j}H)$, $j \in [r]$, are pairwise almost orthogonal. It remains to prove that $H$ is the direct sum of these subspaces. Since $\dim(H) = \sum_{j} \dim(\grm{}{j}H)$, the desired claim follows from Lemma~\ref{lem:orthogonality_vanishing}.
\end{proof}

\section{Original admissible discrete sets}
\setcounter{cdrthm}{4}
\subsubsection{Admissible discrete sets} Let $\highinn{}{\cdot \,, \cdot} \colon H \times H \to \R^r$ be an inner product with values in $\R^r$. Let further $H = H_1 \aplus \cdots \aplus H_r$ be the almost orthogonal decomposition. We write elements $x \in H$ as $x=x_1+\dots+x_r$ with $x_j \in H_j$, $j= 1, \dots, r$.

Consider a subset $S\subset H$. For each $j \in [r]$ and a tuple $z = (z_1, \dots, z_{j-1})$, with $z_i \in H_i$ for $i=1,\dots, j-1$, define the subset $S_{/z} \subseteq S$ as
\begin{equation}\label{eq:SOver}
S_{/z} \coloneqq \left\{ \gamma=\gamma_1 + \dots + \gamma_r \in S \st \gamma_i = z_i \text{ for $i=1, \dots, j-1$} \right\}.
\end{equation}
Moreover, we define the translated set
\begin{equation}\label{eq:Translation}
TS(z) \coloneqq S_{/z} - z_1 - \dots -z_{j-1}.
\end{equation}
We have $TS(z) \subset \filter^j$. In the special case $j=1$, we have the empty vector $z = \varnothing$ and recover
\[
S_{/\varnothing} = TS (\varnothing) = S \subset \filter^1.
\]
\begin{defi}[Admissible discrete sets]
A subset $S \subset H$ is called \emph{admissible discrete} for the inner product $\highinn{}{\cdot\,,\cdot}$, if for all $j \in [r]$ and all tuples $z= (z_1, \dots, z_{j-1})$, with $z_i \in H_i$ for $i=1,\dots, j-1$, the projection $\proj_j(TS(z))$ is a discrete set in $\grm{}{j} H$.%\qedd
\end{defi}

\begin{prop}\label{prop:discrete-discrete}
Every admissible discrete set $S$ in an inner product space $(H, \highinn{}{\cdot \,, \cdot})$ is discrete with respect to the Euclidean topology on $H$.
\end{prop}

\begin{proof}
Assume that this is not the case and let $x$ be an accumulation point of $S$. Then, there exists a sequence $(x^n)_{n \in \N}$ with $ x^n\in S \setminus \{x\}$ for all $n$, which converges to $x$. Proceeding by induction on $j\in [r]$, we show that the equality $x^n_i =x_i$ hold for all $i=1, \dots, j$, and $n \in \N$ large enough. Setting $j=r$, this shows that $x^n = x$ for all large $n \in \N$, contradicting the assumption.

For $j=1$, we use the discreteness of the set $\proj_1(S)$ to infer that the projections $\proj_1(x^n)$ are eventually equal to $\proj_1(x)$. This implies that $x^n_1 = \proj_1^*(\proj_1(x^n)) = \proj_1^*(\proj_1(x)) =x_1$, as required. Assume that the statement holds for $j-1$ and let $z = (x_1, \dots, x_{j-1})$. We use the discreteness of the set $\proj_j(TS(z))$ to infer that for large enough $n$, we have
\[
\proj_j\left(x^n_j\right) = \proj_j\left(x^n - x^n_1-\dots -x^n_{j-1}\right) = \proj_j(x_j).
\]
Applying $\proj_j^*$, we infer that $x^n_j = x_j$ for large enough $n$, as required.
\end{proof}

\setcounter{section}{4}
\section{Tame degenerations}\label{sec:tameness}

Throughout this section, let $H$ be a real vector space of finite dimension endowed with an inner product
\[
\highinn{}{x,y} = \bigl(\highinn{1}{x, y},\highinn{2}{x, y}, \dots, \highinn{r}{x, y} \bigr), \quad x,y \in H,
\]
taking values in $\Lambda = \R^r$ with the lexicographic order $\preceq=\lexeq$. Let
\begin{equation}\label{eq:AlmostOrthoDecompositionState}
H = H_1 \aplus H_2 \aplus \cdots \aplus H_r
\end{equation}
be the almost orthogonal decomposition with $H_j \coloneqq \proj_j^*(\grm{}{j}H)$ (see Lemma~\ref{lem:LiftingLemma}). As before, each element $x\in H$ is written as $x=x_1+ \dots + x_r$ with $x_j \in H_j$, $j \in [r]$.

\begin{conv}[{Reference norm}] \label{con:ref-norm}
In the following, we will need to endow~$H$ with a \emph{reference norm} denoted by $\refnorm{\cdot}$. By equivalence of norms on a given finite dimensional vector space, the discussion in the rest of the paper is independent of the choice of this reference norm. In order to streamline the presentation, we however make the natural choice given by
\begin{equation}\label{eq:ReferenceNorm}
\refnorm{x} = \left(\sum_{j=1}^r \highinn{j}{x_j\,,x_j} \right)^{\frac{1}{2}}, \quad x \in H.%\qedd
\end{equation}%%\qedd
\end{conv}

The fact that the $j^{\rm th}$ component $\highinn{j}{\cdot \,, \cdot}$ of the inner product $\highinn{}{\cdot \,, \cdot}$ is a scalar product on $H_j$ (see Section~\ref{ss:AlmostOrthogonalDecomposition}) ensures that $\refnorm{\cdot}$ is a norm on $H$. Indeed, this property implies that the bilinear form $(x,y) \mapsto \sum_{j} \highinn{j}{x_j\,,y_j}$ is a scalar product on $H$ with associated norm $\refnorm{\cdot}$.

\setcounter{cdrthm}{2}
Let $\innone{t}{\cdot\,, \cdot} \colon H \times H \to \R$, $t\in \R_+$, be scalar products on $H$. Consider a family of vectors $\underline L_t = (L_{t,1}, \dots, L_{t,r}) \in \R_+^r$, $t\in \R_+$.

\begin{defi}[Tame degeneration] \label{def:TameDegenerations}
We say that a family $\innone{t}{\cdot\,, \cdot}$, $t\in \R_+$, of scalar products on $H$ \emph{tamely degenerates} to the inner product $\highinn{}{\cdot\,, \cdot}$ with \emph{parameters} $\underline L_t = (L_{t,1}, \dots, L_{t,r})\in \R_+^r$ if the following conditions hold:
\begin{enumerate}\alphenumi
\item \label{item:tame-a}
For every $j =1, \dots, r-1$, we have
\begin{equation}\label{eq:abstract_tame0}
\lim_{t \to \infty} \frac{L_{t,j}}{L_{t,j+1}} = \infty.
\end{equation}
\item \label{item:tame-b}


For any two distinct indices $i, j \in [r]$, and for every given $\varepsilon>0$, we have
\begin{align}\label{eq:abstract_tame1}
\abs{\innone{t}{x,y}} \le \varepsilon L_{t,\max\{i,j\}}\, \refnorm{x}\cdot\refnorm{y} \quad \forall\,\, x \in H_i, \, y \in H_j,
\end{align}
provided that $t\in \R_+$ is large enough.
\item \label{item:tame-c}
For every $j \in [r]$ and for every given $\delta >0$, we have
\begin{align}\label{eq:abstract_tame2}
\abs{L_{t,j}^{-1} \innone{t}{x, y} - \highinn{j}{x, y}} \le \delta \refnorm{x}\cdot\refnorm{y}, \quad \forall x,y \in H_j,
\end{align}
provided that $t\in \R_+$ is large enough. %\qedd
\end{enumerate}
\end{defi}


In the geometric applications we are interested in, see for example Section~\href{https://ahl.centre-mersenne.org/articles/10.5802/ahl.257/}{10}, it is sometimes necessary to replace~\eqref{item:tame-b} by the following weaker property:
\begin{enumerate}\alphebisnumi
%[label=(\alph*$'$)]
\setcounter{enumi}{1}
\item \label{item:weaktame-b}
There exists a constant $C>0$ such that for all $i, j \in [r]$ and all large $t$, we have
\begin{equation}\label{eq:WeakTameEstimate}
\abs{\innone{t}{x, y} } \le C L_{t, \max\{i,j\}} \refnorm{x}\, \refnorm{y}, \quad \forall\; x \in \filter^i,\, y \in \filter^j.
\end{equation}
\end{enumerate}
For this reason, we introduce the following variant of tame degenerations.

\begin{defi}[$\weak$-tame degeneration]
Notations as above,
let $\innone{t}{\cdot, \cdot}$, $t\in \R_+$, be a family of scalar products on $H$ and $\underline L_t = (L_{t,1}, \dots, L_{t,r}) \in \R_+^r$, $t\in \R_+$, a family of vectors.

We say that the family $\innone{t}{\cdot, \cdot}$, $t\in \R_+$, \emph{degenerates tamely in the weak sense}, or \emph{degenerates $\weak$-tamely}, to $\highinn{}{\cdot\,, \cdot}$ with \emph{parameters} $\underline L_t = (L_{t,1}, \dots, L_{t,r})$ if the properties~\eqref{item:tame-a}, \eqref{item:weaktame-b} and~\eqref{item:tame-c} are satisfied.%\qedd
\end{defi}

\setcounter{cdrthm}{6}
\subsection{Uniform equivalence to orthogonalization in tame degenerations}
 Consider the almost orthogonal decomposition $\aplus_{j=1}^r H_j$ induced by the inner product $\highinn{}{\cdot\,,\cdot}$.
%\colon H \times H \to \R^r$
 As before, each element $x\in H$ is written as $x=x_1+ \dots + x_r$ with $x_j \in H_j$, $j = 1, \dots, r$.


The \emph{orthogonalization} of the inner product $\highinn{}{\cdot\,,\cdot}$ is the new inner product
\begin{align*}
\highinnoneorth{}{\cdot\,,\cdot} &\colon H \times H \to \R^r\\
\highinnoneorth{}{x,y} &\coloneqq \bigl(\highinn{1}{x_1, y_1},\highinn{2}{x_2, y_2}, \dots, \highinn{r}{x_r, y_r} \bigr).
\end{align*}
The following property holds by construction.

\begin{prop}\label{prop:orthogonalization}
Let $\highinnoneorth{}{\cdot\,,\cdot} \colon H \times H \to \R^r$ be the orthogonalization of $\highinn{}{\cdot\,,\cdot}$. Then:
\begin{enumerate}\romanenumi
\item\label{prop5.7.1} The almost orthogonal decomposition induced by $\highinn{}{\cdot\,,\cdot}$ is orthogonal with respect to $\highinnoneorth{}{\cdot\,,\cdot}$. That is, if $x \in H_i$ and $y \in H_j$ with $i \neq j$, then $\highinnoneorth{}{x\,,y} = 0$.
\item\label{prop5.7.2} The reference norm $\refnorm{\cdot}$ coincides with the square root of the pullback $\innoneorth{\underline L}{\cdot\,,\cdot}$ of $\highinnoneorth{}{\cdot\,,\cdot}$ by the vector $\underline L = (1, 1, \dots, 1)$. More precisely,
\[
\refnorm{x} = \sqrt{\innoneorth{\underline L}{x\,,x}}, \quad x \in H.
\]
\item \label{prop:orthogonalization3}
The inner products $\highinn{}{\cdot\,,\cdot}$ and $\highinnoneorth{}{\cdot\,,\cdot}$ have the same tamely degenerating families and the same $\weak$-tamely degenerating families.
\end{enumerate}
\end{prop}

\begin{proof}
Properties~\eqref{prop5.7.1} and~\eqref{prop5.7.2} are clear from the definition of $\highinnoneorth{}{\cdot\,,\cdot}$ and $\refnorm{\cdot}$. Property~\eqref{prop:orthogonalization3} follows from the fact that $\highinn{}{\cdot\,,\cdot}$ and $\highinnoneorth{}{\cdot\,,\cdot}$ induce the same almost orthogonal decomposition of $H$ and $\highinn{j}{x,y} = \highinnoneorth{j}{x,y}$ for all $x,y \in H_j$, see also Theorem~\href{https://ahl.centre-mersenne.org/articles/10.5802/ahl.257/}{11.18}.
\end{proof}

Given a family of vectors $\underline L_t =(L_{1,t}, \dots, L_{r,t}) \in \R_+^r$, $t \in \R_+$, consider the pullback family $\innoneorth{t}{\cdot\,,\cdot} \coloneqq \innoneorth{\underline L_t}{\cdot\,,\cdot}$, $t \in \R_+$ given by
\[
\innoneorth{t}{x,y} = L_{t,1}\highinn{1}{x_1, y_1} + L_{t,2}\highinn{2}{x_2, y_2}+ \dots+L_{r,t}\highinn{r}{x_r, y_r}.
\]
Then, $\innoneorth{t}{\cdot\,, \cdot}$ is a scalar product on $H$ for all $t \in \R_+$. We denote the associated norm on $H$ by
\[
\normorth{x}_t \coloneqq \sqrt{\innoneorth{t}{x, x}} = \left(\sum_{j=1}^r L_{t,j} \qf_j(x_j) \right)^{1/2}, \quad x \in H,
\]
where $\qf_j$ is the $j^{\rm th}$ component of the quadratic form $\qf$ associated to $\highinn{}{\cdot\,,\cdot}$, and we have
\[
\qf_j(z) = \highinn{j}{z, z}, \quad z \in H_j.
\]
\begin{prop}[Uniform equivalence to orthogonalization pullbacks] \label{prop:apriori_estimate}
Consider an inner product $\highinn{}{\cdot\,, \cdot} \colon H \times H \to \R^r$. Let $\innone{t}{\cdot \,, \cdot} \colon H \times H \to \R$, $t \in \R_+$, be a family of scalar products which $\weak$-tamely degenerates to $\highinn{}{\cdot\,, \cdot}$ with parameters $\underline L_t \in \R_+^r$. Denote by $\norm{x}_t \coloneqq \sqrt{\innone{t}{x,x}}$, $t \in \R_+$, the associated norms on $H$.

Then, the two families of norms $\norm{\cdot}_t $ and $\normorth{\cdot}_t$ are uniformly equivalent for large~$t$. That is, there exists a constant $D >0$ such that
\[
D^{-1/2} \normorth{\cdot}_t \le \norm{\cdot}_t \le D^{1/2} \normorth{\cdot}_t,
\]
provided that $t \in \R_+$ is large enough. Equivalently,
\begin{equation}\label{eq:apriori_estimate}
D^{-1} \sum_{j=1}^r L_{t,j} \left\|x_j\right\|^2_j \le \normbsq{x}_t \le D \sum_{j=1}^r L_{t,j} \normbsq{x_j}_j \quad \forall x\in H,
\end{equation}
provided that $t \in \R_+$ is large enough.
\end{prop}

\begin{proof}
Writing $x \in H$ as $x = x_1+ \dots + x_r$, with $x_j\in H_j$ for $j=1, \dots, r$, we have
\[
\normsq{x}_t = \sum_{j=1}^r \normbsq{x_j}_t + 2 \sum_{i <j} \innone{t}{x_i, x_j}.
\]
It follows from~\eqref{eq:abstract_tame2} that there exists a uniform constant $D' >0$ such that
\begin{multline*}
\frac 1{D'} \sum_{j=1}^r L_{t,j} \normbsq{x_j}_j \ge \sum_{j=1}^r \normbsq{x_j}_t \ge D' \sum_{j=1}^r L_{t,j} \normbsq{x_j}_j \\
\forall \, x\in H \;\text{and}\; \forall\, t\in \R_+ \text{ large enough}.
\end{multline*}
Fix now a pair of indices $i,j\in [r]$ with $i<j$. Combining the inequality
\begin{equation}\label{eq:trick}
2 ab \le\frac{1}{\delta} a^2 + \delta b^2, \quad a,b \ge 0, \, \delta >0,
\end{equation}
with assumption~\eqref{eq:WeakTameEstimate}, we conclude that for every $\delta >0$, there exists a constant $C_\delta >0$ with
\begin{multline*}
\abs{\innone{t}{x_i, x_j}} \le C L_{t,j} \norm{x_i}_i \cdot \norm{x_j}_j \le L_{t,j}\left(C_\delta \normbsq{x_i}_i + \delta \normbsq{x_j}_j\right) \\ 
\forall \, x\in H \quad\text{and}\quad t \in \R_+ \text{ large enough}.
\end{multline*}
Taking into account that $L_{t,i} /L_{t,j} \to \infty$ for $t \to \infty$, it follows that
\[
\abs{\innone{t}{x_i, x_j}} \le \delta L_{t,i} \normbsq{x_i}_i + \delta L_{t,j} \normsq{x_j}_j
\]
for $t$ large. Choosing $\delta >0$ small and combining the above estimates, we conclude.
\end{proof}

\section{A finiteness lemma} \label{sec:FinitenessLemma}
In this section we state and prove the following finiteness lemma. It is a key ingredient in the arguments of the next sections.

Consider an inner product $\highinn{}{\cdot\,,\cdot} \colon H\times H \to \R^r$. Let $\innone{t}{\cdot\,,\cdot} \colon H\times H \to \R$, $t\in \R_+$, be scalar products which $\weak$-tamely degenerate to $\highinn{}{\cdot\,,\cdot}$ with parameters $\underline L_t \in \R_+^r$. Denote by $\norm{x}_t^2 = \innone{t}{x,x}$, $t \in \R_+$, the associated norms on $H$.

\begin{lemm}[Finiteness Lemma] \label{lem:finiteness-discrete}
Let $S \subset H$ be an admissible discrete set for $\highinn{}{\cdot\,,\cdot}$. Then, the following holds. For any compact subset $B \subseteq H$, there exists a finite subset $\S_B \subset S$ such that for any large enough $t$, we have
\[
\min_{\gamma \in S} \norm{x -\gamma}_t = \min_{\gamma \in \S_B} \norm{x -\gamma}_t \quad \forall \: x \in B.
\]
\end{lemm}

For scalar products (i.e., when $r=1$), this lemma trivially follows from the compactness of $B$. On the contrary, the proof for higher values of $r$ is quite involved.

As before, we use the almost orthogonal decomposition $H = \aplus_{j=1}^r H_j$ and decompose each $x \in H$ as $x = x_1 + \dots +x_r$ with $x_j \in H_j$, $j\in[r]$. Moreover by our Convention~\ref{con:ref-norm}, we endow $H$ with the reference norm
\[
\refnorm{x} = \left(\sum_{j=1}^r \highinn{j}{x_j\,,x_j} \right)^{\frac{1}{2}}, \quad x \in H.
\]

\begin{proof}
The construction of the set $\S_B$ will be by induction. We will successively choose appropriate radii $R_1, \dots, R_r \in \R_+$ such that the sets
\[
S^0 \coloneqq S\supseteq S^1 \supseteq S^2 \supseteq S^3 \supseteq \dots \supseteq S^{r}
\]
defined by
\begin{equation}\label{eq:define-successive-sets}
S^j = S^j(R_1, \dots, R_j) \coloneqq \bigl\{ \gamma = \gamma_1 + \dots +\gamma_r \in S \bigm|  \refnorm{\gamma_k} \le R_k \quad\text{for all}\; k \le j \bigr\}
\end{equation}
satisfy the property
\begin{equation}\label{eq:finiteness-wish}
\min_{\gamma \in S} \normsq{x -\gamma}_t = \min_{\gamma \in S^j} \normsq{x -\gamma}_t \quad \forall \: x\in B,
\end{equation}
provided that $t$ is large enough. Moreover, we define $Z^j=Z^j(R_1, \dots, R_j)$ to be the set of all tuples $(\gamma_1, \dots, \gamma_j)$ appearing in the decomposition $\gamma =\gamma_1 + \dots +\gamma_r$ for $\gamma \in S^j$, and show that
\begin{equation}\label{eq:finiteness-wish2}
Z^j \text{ is a finite set.}
\end{equation}
The last set $S^r$ coincides with $Z^r$, which is finite. The claim then follows by setting $\S_B \coloneqq S^r$.


The properties~\eqref{eq:finiteness-wish} and~\eqref{eq:finiteness-wish2} are trivially satisfied for $S^0 = S$ and $Z^0 =\emptyset$.

Suppose we have already constructed radii $R_1, \dots, R_j \in \R_+$ such that~\eqref{eq:finiteness-wish} holds for $S^j$ defined by~\eqref{eq:define-successive-sets}, and~\eqref{eq:finiteness-wish2} is also satisfied for $Z^j$.


For $\gamma \in S^j$ and $x \in H$, we then write
\begin{multline*}
\normsq{x - \gamma}_t = \normbsq{x-(\gamma_1 + \dots +\gamma_j)}_t + \normbsq{\gamma_{j+1} + \dots + \gamma_r}_t\\
- 2\sum_{k = j+1}^r \innone{t}{x- (\gamma_1 + \dots + \gamma_j), \gamma_k}
\end{multline*}
for the decomposition $\gamma = \gamma_1 + \dots + \gamma_r$, $\gamma_j\in H_j$. Combing~\eqref{eq:WeakTameEstimate} and the inequality~\eqref{eq:trick}, we conclude that for every $\delta >0$, there exists a constant $C_\delta >0$ such that
\[
\sum_{k = j+1}^r \abss{\innone{t}{x- (\gamma_1 + \dots + \gamma_j), \gamma_k}} \le \sum_{k =j+1}^r L_{t,k} \left(C_\delta \refnormsq{x- (\gamma_1 + \dots + \gamma_j)} + \delta \refnormsq{\gamma_k}\right)
\]
holds true for all $x\in B$ and $\gamma \in S^j$. By the definition of $S^j = S^j(R_1, \dots, R_j)$, and the compactness of $B$, the elements $x - (\gamma_1+ \dots+\gamma_j)$, for $\gamma \in S^j$ and $x \in B$,\pagebreak{} are all contained in a compact subset $\tilde{B}$ of $H$. In particular, using~\eqref{eq:abstract_tame0}, there exists a constant $D_\delta >0$ such that
\begin{multline*}
2 \sum_{k = j+1}^r \abss{\innone{t}{x- (\gamma_1 + \dots + \gamma_j), \gamma_k}}\\
\le D_\delta L_{t,j+1} + \sum_{k =j+1}^r L_{t,k} \delta \refnormsq{\gamma_k}
\forall \: \gamma \in S^j \;\text{ and}\quad \forall \: x\in B.
\end{multline*}
Applying Proposition~\ref{prop:apriori_estimate} and choosing $\delta >0$ small, we get the estimate
\begin{multline}\label{eq:estimate-induction1}
\normsq{x - \gamma}_t - \normbsq{x-(\gamma_1 + \dots +\gamma_j)}_t\\
\begin{aligned}
&\ge L_{t,j+1} \left(D \refnormbsq{\gamma_{j+1}} - \delta \refnormbsq{\gamma_{j+1}} - D_\delta \right) + \sum_{k=j+2}^r L_{t,k} (D - \delta) \refnormsq{\gamma_k} \\
& \ge L_{t,j+1} \left(\frac{D}{2} \refnormbsq{\gamma_{j+1}} - D_\delta\right),
\end{aligned}
\end{multline}
holding uniformly for all $x \in B$, $\gamma \in S^j$, and $t$ large.

To proceed, we choose a large positive real $D' \in\R_+$, to be determined below, and choose $R_{j+1} \in \R_+$ large enough so that $\frac{D}{2} R_{j+1}^2 - D_\delta >3D'$. We then define the set $S^{j+1} = S^j(R_1, \dots, R_{j+1}) \subseteq S$ by~\eqref{eq:define-successive-sets}.



We claim that, for large $t$, and all $\gamma \in S^j \setminus S^{j+1}$ and $x \in B$, there exists $\tilde \gamma \in S^{j+1}$ such that we have the estimate
\begin{equation}\label{eq:FinitenessFinalStep}
\normsq{x - \gamma}_t > \normbsq{x - \tilde \gamma}_t.
\end{equation}
This will prove the desired claim~\eqref{eq:finiteness-wish}.



Indeed for all $\gamma \in S^j \setminus S^{j+1}$, we have $\refnormsq{\gamma_{j+1}} \ge R_{j+1}^2$. The estimate~\eqref{eq:estimate-induction1} and the choice of $R_{j+1}$ gives
\begin{align}\label{eq:estimate-induction2}
\normsq{x - \gamma}_t & > \normsq{x-(\gamma_1 + \dots +\gamma_j)}_t + 3D' L_{t, j+1}
\end{align}
for large enough $t$.



For all $z = (\gamma_1, \dots, \gamma_j) \in Z^j$, set
\[
S_{/z} \coloneqq \bigl\{ \theta =\theta_1+ \dots+\theta_r \in S\bigm| \theta_1=\gamma_1, \dots, \theta_j = \gamma_j\bigr\}.
\]

Since $B$ is compact and $Z^j$ is finite, the union $\bigcup_{z\in Z^j} \left(B -z\right)$ is compact. Using the properties~\eqref{eq:abstract_tame0}--\eqref{eq:abstract_tame1}--\eqref{eq:abstract_tame2}--\eqref{eq:WeakTameEstimate}, we choose $D'$ large enough so that for all $z =(\gamma_1, \dots, \gamma_j)\in Z^j$, we can find an element $\tilde \gamma = \tilde\gamma_z \in S_{/z}$ such that
\begin{align}
\normbsq{\tilde \gamma_{j+1} + \dots + \tilde \gamma_r}_t &\leq D' L_{t, j+1} \label{eq:estimate-induction3} \\
\intertext{and}
\abss{\innone{t}{x-\gamma_1-\dots-\gamma_j, \tilde \gamma_{j+1} + \dots + \tilde \gamma_r }} &\leq D' L_{t,j+1} \quad \text{for all $x\in B$},\label{eq:estimate-induction4}
\end{align}
provided that $t$ is large enough. Choosing $R_{j+1}$ large enough, we can ensure that $\tilde \gamma = \tilde\gamma_z$ belongs to $S^{j+1}$ for all $z \in Z^j$.

To conclude, we combine~\eqref{eq:estimate-induction2}, \eqref{eq:estimate-induction3}, and~\eqref{eq:estimate-induction4} to get for all $\gamma \in S^j \setminus S^{j+1}$ with $\gamma \in S_{/z}$,
\begin{align*}
\normsq{x - \gamma}_t & > \normsq{x-(\gamma_1 + \dots +\gamma_j)}_t + 3D' L_{t, j+1} \\
&\ge \normsq{x-(\gamma_1 + \dots +\gamma_j)}_t + \normbsq{\tilde \gamma_{j+1} + \dots + \tilde \gamma_r}_t + 2 D' L_{t,j+1} \\
& = \normbsq{x- \tilde \gamma}_t - 2 \binnone{t}{x-(\gamma_1 + \dots +\gamma_j), \tilde \gamma_{j+1} + \dots + \tilde \gamma_r } + 2 D' L_{t,j+1} \\
&\ge \normsq{x- \tilde \gamma}_t.
\end{align*}

This establishes~\eqref{eq:FinitenessFinalStep}.



It remains to show that the set $Z^{j+1} =Z^{j+1}(R_1, \dots, R_{j+1})$ is finite. This will be a direct consequence of the finiteness of $Z^j$ and the fact that $S \subset H$ is admissible discrete. Fix $z =(z_1, \dots, z_j)$ in $Z^j$ and consider the subset $S_{/z}\!\!\subset\!S$. Since $S$
%\subset\!H$
is admissible discrete, the subset $\{\gamma_{j+1} \gamma = \gamma_1 + \dots + \gamma_r \in S_{/z}\} \subset H_{j+1}$ is a discrete subset of $H_{j+1} \cong \grm{}{j+1}H$. Restricting to $S^{j+1} \cap S_{/z}$, we obtain that $\{\gamma_{j+1} \gamma = \gamma_1 + \dots + \gamma_r \in S_{/z} \cap S^{j+1}\}$ is finite. By finiteness of $Z^j$, we conclude that $Z^{j+1}$ is finite.
\end{proof}

\begin{thebibliography}{99}
\bibitem[Bir40]{Bir-Lattice} G.Birkhoff, \emph{Lattice theory}, Colloquium Publications25, American Mathematical Society,1940.
\bibitem[HW52]{HW52} M.Hausner and J.G.Wendel, ``Ordered vector spaces'', \emph{Proc. Am. Math. Soc.}\textbf{3}(1952),977--982.
\end{thebibliography}
\end{document}
